\documentclass[border=4pt]{standalone}
\usepackage{fontspec}
\setmainfont{Times New Roman}
\usepackage[russian]{babel}
\usepackage{tikz}
\usetikzlibrary{arrows.meta,positioning,calc,decorations.pathreplacing,fit,backgrounds}
% палитра (dataviz, светлая тема)
\definecolor{cblue}{HTML}{2A78D6}\definecolor{corange}{HTML}{EB6834}\definecolor{caqua}{HTML}{1BAF7A}
\definecolor{cyellow}{HTML}{EDA100}\definecolor{ink}{HTML}{0B0B0B}\definecolor{ink2}{HTML}{52514E}
\definecolor{grid}{HTML}{E4E3DF}\definecolor{paper}{HTML}{FCFCFB}
\tikzset{
  >={Stealth[length=2mm]},
  blk/.style={draw=ink2, rounded corners=2pt, minimum height=6mm, minimum width=11mm, font=\footnotesize, fill=white},
  orig/.style={blk, fill=grid!60},
  repl/.style={blk, draw=corange, fill=corange!12, line width=0.7pt},
  plus/.style={circle, draw=ink2, inner sep=0.6pt, font=\scriptsize, fill=white},
  note/.style={font=\scriptsize, text=ink2, align=center},
}

\begin{document}
\begin{tikzpicture}[x=1cm,y=1cm]
  \def\N{12}\def\s{0.22}
  % полное каузальное внимание
  \begin{scope}[shift={(0,0)}]
    \node[font=\small\bfseries, anchor=south west] at (0,\N*\s+0.15) {Внимание: исходное};
    \foreach \t in {0,...,11}{\foreach \j in {0,...,\t}{
      \fill[cblue!55] (\j*\s,{(\N-1-\t)*\s}) rectangle ++(\s*0.9,\s*0.9);}}
    \node[note, anchor=north] at (\N*\s/2,-0.1) {каждый токен (строка) смотрит\\на все предыдущие (столбцы)};
  \end{scope}
  % окно + сток
  \begin{scope}[shift={(3.6,0)}]
    \node[font=\small\bfseries, anchor=south west] at (0,\N*\s+0.15) {Замена: окно + сток};
    \foreach \t in {0,...,11}{\foreach \j in {0,...,\t}{
      \pgfmathtruncatemacro{\d}{\t-\j}
      \ifnum\d<3 \fill[corange!70] (\j*\s,{(\N-1-\t)*\s}) rectangle ++(\s*0.9,\s*0.9);
      \else\ifnum\j=0 \fill[corange!70] (\j*\s,{(\N-1-\t)*\s}) rectangle ++(\s*0.9,\s*0.9);
      \else \fill[grid] (\j*\s,{(\N-1-\t)*\s}) rectangle ++(\s*0.9,\s*0.9);\fi\fi}}
    \node[note, anchor=north] at (\N*\s/2,-0.1) {только ближайшие токены\\и первый (сток)};
  \end{scope}
  % нейроны MLP
  \begin{scope}[shift={(7.4,0.4)}]
    \node[font=\small\bfseries, anchor=south west] at (0,2.35) {MLP: какие нейроны считать};
    \foreach \r/\lab/\pat in {1.6/{исходный}/{1,1,1,1,1,1,1,1,1,1,1,1,1,1,1,1}, 0.6/{замена}/{0,1,0,0,1,0,0,0,1,0,0,1,0,0,0,0}}{
      \node[note, anchor=east] at (-0.1,\r+0.12) {\lab};
      \foreach \v [count=\c from 0] in \pat {
        \ifnum\v=1 \fill[\ifdim\r pt>1pt cblue!55\else corange!70\fi] (\c*0.24,\r) rectangle ++(0.2,0.24);
        \else \fill[grid] (\c*0.24,\r) rectangle ++(0.2,0.24);\fi}}
    \node[note, anchor=north west, align=left] at (0,0.45) {для каждого токена предиктор выбирает\\25\% нейронов; остаток — линейная поправка};
  \end{scope}
\end{tikzpicture}
\end{document}
